% ============================================================
%  CHAPTER 7: Comparative Analysis
% ============================================================

\chapter{Comparative Analysis}
\label{chap:comparison}

The preceding chapters examined KZG (Chapter~\ref{chap:kzg}),
Multilinear KZG (Chapter~\ref{chap:mkzg}),
Bulletproofs (Chapter~\ref{chap:bulletproofs}), and
Dory (Chapter~\ref{chap:dory}) each in isolation. This chapter brings
all four schemes together and analyses them side by side across every
benchmarked metric: setup time, commitment time, proving time,
verification time, proof size, and the impact of
Pippenger's multi-scalar multiplication optimisation. All benchmark
results were collected on the same machine under identical conditions
using Rust and the \texttt{arkworks} library, with sizes ranging from
$n = 4$ to $n = 1024$ field elements and the three curves
BLS12-381, BN254, and BLS12-377.

A key point before comparing: the four schemes do not commit to
identical mathematical objects. KZG commits to a univariate polynomial
of degree $d$; MKZG commits to a multilinear polynomial in $\mu$
variables; Bulletproofs and Dory commit to a vector of $n$ scalars
(equivalently, a multilinear polynomial in $\log_2 n$ variables).
Throughout this chapter, ``size $n$'' means: the degree bound $d = n$
for KZG, the number of variables $\mu = \log_2 n$ (poly size $2^\mu = n$)
for MKZG, and the vector length $n$ for Bulletproofs and Dory. In all
cases the prover holds approximately $n$ field elements, making the
comparison on equal footing with respect to the committed data volume.



% ============================================================
\section{Setup Time}
\label{sec:setup-comparison}
% ============================================================

Table~\ref{tab:setup-all} shows setup times across all four schemes.
KZG computes $n$ powers of $\tau$ in $\mathbb{G}_1$ and $\mathbb{G}_2$;
MKZG evaluates the multilinear Lagrange basis at a secret point;
Bulletproofs hashes $n$ random curve points; and Dory generates
random generators and precomputes $O(\log n)$ inner pairing products.

\begin{table}[H]
\centering
\caption{Setup time (ms) at representative sizes. KZG: degree $n$;
MKZG: $\mu = \log_2 n$ variables; Bulletproofs/Dory: vector length $n$.}
\label{tab:setup-all}
\adjustbox{max width=\textwidth}{%
\begin{tabular}{rlrrrrrrrrrrrr}
\toprule
& & \multicolumn{3}{c}{\textbf{KZG}}
  & \multicolumn{3}{c}{\textbf{MKZG}}
  & \multicolumn{3}{c}{\textbf{Bulletproofs}}
  & \multicolumn{3}{c}{\textbf{Dory}} \\
\cmidrule(lr){3-5}\cmidrule(lr){6-8}\cmidrule(lr){9-11}\cmidrule(lr){12-14}
$n$ & & \textbf{381} & \textbf{254} & \textbf{377}
    & \textbf{381} & \textbf{254} & \textbf{377}
    & \textbf{381} & \textbf{254} & \textbf{377}
    & \textbf{381} & \textbf{254} & \textbf{377} \\
\midrule
4    & & 6.10   & 2.46  & 5.70   & 3.10  & 1.37  & 2.75  & 0.65  & 0.09  & 0.62  & 14.13 & 7.08  & 14.43 \\
16   & & 18.31  & 10.49 & 19.03  & 7.46  & 3.80  & 7.06  & 2.19  & 0.38  & 2.35  & 20.89 & 11.44 & 25.15 \\
64   & & 69.46  & 32.35 & 76.58  & 19.05 & 10.41 & 19.62 & 8.05  & 1.22  & 9.67  & 37.90 & 19.34 & 42.89 \\
256  & & 240.85 & 129.52 & 296.95 & 61.34 & 35.10 & 65.08 & 37.49 & 5.46  & 37.06 & 63.80 & 33.50 & 73.61 \\
1024 & & 1027.47 & 544.87 & 1156.37 & 235.94 & 129.63 & 239.44 & 154.91 & 18.36 & 149.98 & 110.96 & 57.59 & 126.04 \\
\bottomrule
\end{tabular}}%
\end{table}

\textbf{KZG has the most expensive setup.} At $n = 1024$ on BLS12-381,
KZG takes 1027\,ms, which is $4.4\times$ more than MKZG (236\,ms),
$6.6\times$ more than Bulletproofs (155\,ms), and $9.3\times$ more
than Dory (111\,ms). KZG must compute $n+1$ structured powers
$[\tau^i]_1$ with all distinct scalars, whereas MKZG uses $n$
fixed-base scalar multiplications on the same generator $g_1$,
allowing window-table precomputation to be reused across all $n$
multiplications. This fixed-base advantage accounts for MKZG's
$4.4\times$ lower constant despite both being $O(n)$.
Bulletproofs setup (18\,ms on BN254 at $n = 1024$) is negligible,
consisting only of $n$ random point hashes; Dory is cheapest among
pairing schemes as it uses unstructured generators.

\textbf{MKZG setup is exponential in the number of variables.}
The table rows use poly\-size $n = 2^{\mu}$, so setup appears linear
in $n$, but it is \emph{exponential in $\mu$}: each extra variable
doubles the SRS size and cost. The rows $n = 4, 16, 64, 256, 1024$
correspond to $\mu = 2, 4, 6, 8, 10$ variables; eight more variables
grows the SRS $256\times$ (from 4 to 1024 elements) and setup from
3.10\,ms to 235.94\,ms on BLS12-381. Extrapolating: $\mu = 20$
requires $2^{20} \approx 10^6$ group elements ($\approx 242$\,GB),
and $\mu = 30$ exceeds 240\,TB, making unoptimised MKZG impractical
beyond $\mu \approx 20$--$24$.

\textbf{BN254 is consistently fastest.} BN254 is $\approx 1.9\times$
faster than BLS12-381 for KZG/MKZG/Dory, and over $8\times$ faster
for Bulletproofs setup (where cost is dominated by point hashing,
which scales with byte length rather than field arithmetic).

% ============================================================
\section{Commitment Time}
\label{sec:commit-comparison}
% ============================================================

The commitment operation in all four schemes reduces to a
multi-scalar multiplication (MSM) of size $n$. Consequently,
Pippenger's algorithm has the most direct and uniform impact here.


\begin{table}[H]
\centering
\caption{Commit time (ms) at representative sizes with Pippenger.}
\label{tab:commit-all}
\adjustbox{max width=\textwidth}{%
\begin{tabular}{rlrrrrrrrrrrrr}
\toprule
& & \multicolumn{3}{c}{\textbf{KZG}}
  & \multicolumn{3}{c}{\textbf{MKZG}}
  & \multicolumn{3}{c}{\textbf{Bulletproofs}}
  & \multicolumn{3}{c}{\textbf{Dory}} \\
\cmidrule(lr){3-5}\cmidrule(lr){6-8}\cmidrule(lr){9-11}\cmidrule(lr){12-14}
$n$ & & \textbf{381} & \textbf{254} & \textbf{377}
    & \textbf{381} & \textbf{254} & \textbf{377}
    & \textbf{381} & \textbf{254} & \textbf{377}
    & \textbf{381} & \textbf{254} & \textbf{377} \\
\midrule
4    & & 0.83  & 0.39  & 0.74  & 0.68  & 0.36  & 0.58  & 0.76  & 0.40  & 0.63  & 2.72  & 1.93  & 3.29  \\
16   & & 1.68  & 1.46  & 1.68  & 1.33  & 0.80  & 1.38  & 1.42  & 0.70  & 1.62  & 6.11  & 3.16  & 6.03  \\
64   & & 4.32  & 2.12  & 4.69  & 4.10  & 2.27  & 4.13  & 3.92  & 1.86  & 4.34  & 13.39 & 6.71  & 13.66 \\
256  & & 12.86 & 6.52  & 13.94 & 12.04 & 6.21  & 12.12 & 12.98 & 5.25  & 13.68 & 33.66 & 17.16 & 34.27 \\
1024 & & 40.37 & 20.88 & 42.00 & 39.48 & 20.06 & 39.05 & 39.05 & 15.41 & 39.98 & 90.96 & 49.25 & 101.81 \\
\bottomrule
\end{tabular}}%
\end{table}

\textbf{KZG, MKZG, and Bulletproofs commit in nearly identical time.}
At $n = 1024$ on BLS12-381, all three take 39--41\,ms with Pippenger
enabled. This is expected: all three compute a single MSM of $n$
elements in $\mathbb{G}_1$, so the commit cost is determined purely by
the MSM cost. The near-identical values directly validate that
Pippenger's algorithm dominates and the scheme-specific overhead is
negligible.

\textbf{Dory commit is roughly $2.3\times$ more expensive.} At
$n = 1024$ on BLS12-381, Dory takes 91\,ms compared to $\approx$40\,ms
for the other three. Dory's commitment involves $2^\nu$ row MSMs of
size $2^\sigma$ each (total $O(n)$ scalar multiplications), followed by
a second-tier MSM and a pairing. The additional $\mathbb{G}_T$ pairing
at the end accounts for most of the overhead vs.\ the pure
$\mathbb{G}_1$ MSM of the other schemes.

\textbf{The BN254 advantage is $\approx 1.9\times$ for most schemes.}
Across KZG, MKZG, and Bulletproofs at $n = 1024$, BN254 is
$1.93\times$, $1.97\times$, and $2.53\times$ faster than BLS12-381
respectively. Bulletproofs shows a larger BN254 advantage because its
commit implementation uses integer-only point hashing paths that
benefit more from BN254's smaller field size.

% ============================================================
\section{Proving Time}
\label{sec:prove-comparison}
% ============================================================

Proving time is where the four schemes diverge most.
Tables~\ref{tab:prove-all} and~\ref{tab:prove-ratio} report raw times
and ratios relative to KZG on BLS12-381.

\begin{table}[H]
\centering
\caption{Prove time (ms) at representative sizes.}
\label{tab:prove-all}
\adjustbox{max width=\textwidth}{%
\begin{tabular}{rlrrrrrrrrrrrr}
\toprule
& & \multicolumn{3}{c}{\textbf{KZG (Pip)}}
  & \multicolumn{3}{c}{\textbf{MKZG}}
  & \multicolumn{3}{c}{\textbf{Bulletproofs (Pip)}}
  & \multicolumn{3}{c}{\textbf{Dory}} \\
\cmidrule(lr){3-5}\cmidrule(lr){6-8}\cmidrule(lr){9-11}\cmidrule(lr){12-14}
$n$ & & \textbf{381} & \textbf{254} & \textbf{377}
    & \textbf{381} & \textbf{254} & \textbf{377}
    & \textbf{381} & \textbf{254} & \textbf{377}
    & \textbf{381} & \textbf{254} & \textbf{377} \\
\midrule
4    & & 0.67  & 0.31  & 0.76  & 1.07  & 0.67  & 0.91  & 3.34  & 1.70  & 3.01  & 22.92 & 15.36 & 27.74 \\
16   & & 1.66  & 1.52  & 1.62  & 4.63  & 2.43  & 4.70  & 10.66 & 6.38  & 11.77 & 45.77 & 28.64 & 55.45 \\
64   & & 4.46  & 2.30  & 5.23  & 17.61 & 9.49  & 18.43 & 38.90 & 21.81 & 43.59 & 89.20 & 51.62 & 104.92 \\
256  & & 12.19 & 6.84  & 15.35 & 72.18 & 36.78 & 71.17 & 142.95 & 84.38 & 158.03 & 219.85 & 92.87 & 182.57 \\
1024 & & 42.10 & 23.51 & 46.16 & 263.37 & 149.15 & 272.35 & 650.35 & 415.67 & 582.06 & 264.31 & 173.06 & 309.89 \\
\bottomrule
\end{tabular}}%
\end{table}

\begin{table}[H]
\centering
\caption{Prove time ratio relative to KZG on BLS12-381 at each size.
A ratio of $k$ means the scheme takes $k\times$ longer to prove than KZG.}
\label{tab:prove-ratio}
\begin{tabular}{rllll}
\toprule
$n$ & \textbf{KZG} & \textbf{MKZG} & \textbf{Bulletproofs} & \textbf{Dory} \\
\midrule
4    & $1\times$  & $1.6\times$  & $5.0\times$   & $34.3\times$ \\
16   & $1\times$  & $2.8\times$  & $6.4\times$   & $27.6\times$ \\
64   & $1\times$  & $3.9\times$  & $8.7\times$   & $20.0\times$ \\
256  & $1\times$  & $5.9\times$  & $11.7\times$  & $18.0\times$ \\
1024 & $1\times$  & $6.3\times$  & $15.4\times$  & $6.3\times$  \\
\bottomrule
\end{tabular}
\end{table}

\textbf{KZG} is fastest: a single $O(n)$ MSM for the quotient
polynomial gives 42\,ms at $n = 1024$ on BLS12-381.

\textbf{MKZG and Dory} both land near 264\,ms at $n = 1024$
($6.3\times$ KZG). MKZG performs $\mu$ MSMs of halving size ($O(n\log n)$
total); Dory runs $\sigma = 5$ fold rounds of row MSMs, also $O(n\log n)$
in practice. Dory starts at $34\times$ KZG at $n = 4$ (fixed round
overhead dominates) but converges to $6\times$ as $n$ grows and the
$O(n)$ row MSMs become the dominant cost.

\textbf{Bulletproofs} is the slowest prover: 650\,ms at $n = 1024$
($15.4\times$ KZG). Its $\log_2 n = 10$ halving rounds each require
two MSMs, totalling $2n - 2$ scalar multiplications with no shared
structure to batch.

% ============================================================
\section{Verification Time}
\label{sec:verify-comparison}
% ============================================================

Tables~\ref{tab:verify-all} and~\ref{tab:verify-ratio} show
verification times and their ratio relative to KZG.

\begin{table}[H]
\centering
\caption{Verify time (ms) at representative sizes.}
\label{tab:verify-all}
\adjustbox{max width=\textwidth}{%
\begin{tabular}{rlrrrrrrrrrrrr}
\toprule
& & \multicolumn{3}{c}{\textbf{KZG (Pip)}}
  & \multicolumn{3}{c}{\textbf{MKZG}}
  & \multicolumn{3}{c}{\textbf{Bulletproofs (Pip)}}
  & \multicolumn{3}{c}{\textbf{Dory}} \\
\cmidrule(lr){3-5}\cmidrule(lr){6-8}\cmidrule(lr){9-11}\cmidrule(lr){12-14}
$n$ & & \textbf{381} & \textbf{254} & \textbf{377}
    & \textbf{381} & \textbf{254} & \textbf{377}
    & \textbf{381} & \textbf{254} & \textbf{377}
    & \textbf{381} & \textbf{254} & \textbf{377} \\
\midrule
4    & & 4.75  & 2.65  & 4.82  & 6.41  & 4.85  & 6.79  & 2.53  & 1.46  & 2.73  & 29.35 & 24.45 & 33.11 \\
16   & & 4.29  & 3.90  & 4.74  & 10.35 & 6.96  & 11.61 & 8.37  & 5.59  & 9.14  & 45.38 & 29.70 & 56.15 \\
64   & & 4.01  & 2.77  & 7.85  & 13.88 & 9.76  & 16.65 & 29.77 & 17.10 & 34.16 & 71.09 & 42.03 & 76.57 \\
256  & & 3.58  & 2.53  & 4.38  & 17.97 & 12.37 & 21.42 & 116.46 & 64.82 & 120.30 & 81.63 & 54.69 & 92.02 \\
1024 & & 3.81  & 2.53  & 4.65  & 21.85 & 15.98 & 26.61 & 496.26 & 329.64 & 499.11 & 95.82 & 61.90 & 112.07 \\
\bottomrule
\end{tabular}}%
\end{table}

\begin{table}[H]
\centering
\caption{Verify time ratio relative to KZG (Pip) on BLS12-381 at each size.}
\label{tab:verify-ratio}
\begin{tabular}{rllll}
\toprule
$n$ & \textbf{KZG (Pip)} & \textbf{MKZG} & \textbf{Bulletproofs (Pip)} & \textbf{Dory} \\
\midrule
4    & $1\times$ & $1.3\times$  & $0.5\times$  & $6.2\times$  \\
16   & $1\times$ & $2.4\times$  & $2.0\times$  & $10.6\times$ \\
64   & $1\times$ & $3.5\times$  & $7.4\times$  & $17.7\times$ \\
256  & $1\times$ & $5.0\times$  & $32.5\times$ & $22.8\times$ \\
1024 & $1\times$ & $5.7\times$  & $130.3\times$ & $25.1\times$ \\
\bottomrule
\end{tabular}
\end{table}

\textbf{KZG} verifies in exactly two pairings: 3.6--4.8\,ms on
BLS12-381, flat across all $n$. This $O(1)$ cost is its defining
advantage.

\textbf{MKZG} requires $\mu + 1$ pairings ($O(\log n)$). At
$\approx 1.90$\,ms per pairing on BLS12-381, the expected costs are
$3 \times 1.90 = 5.70$\,ms at $n = 4$ (observed: 6.41\,ms) and
$11 \times 1.90 = 20.9$\,ms at $n = 1024$ (observed: 21.85\,ms).
Even at $n = 16384$ ($\mu = 14$) verify takes only $\approx$33\,ms.

\textbf{Dory} is also $O(\log n)$ but with a higher constant: 95.8\,ms
at $n = 1024$ on BLS12-381 ($25\times$ KZG). Each of the $\sigma = 5$
fold rounds involves $\mathbb{G}_T$ arithmetic (six 576-byte elements
per round), which is costlier than MKZG's $\mathbb{G}_2$ pairings.

\textbf{Bulletproofs} verify is $O(n)$ and impractical at scale:
496\,ms at $n = 1024$ on BLS12-381 ($130\times$ KZG). Verify grows
$\approx 196\times$ from $n = 4$ to $n = 1024$, consistent with
linear scaling.


% ============================================================
\section{Proof Size}
\label{sec:proofsize-comparison}
% ============================================================

Table~\ref{tab:proofsize-all} gives proof sizes across all schemes.
The four schemes span a $433\times$ range at $n = 1024$ on BLS12-381:
\[
\text{KZG}\ (48\,\text{B})\ \ll\ \text{MKZG}\ (480\,\text{B})\ <\
\text{Bulletproofs}\ (2{,}216\,\text{B})\ \ll\ \text{Dory}\ (20{,}784\,\text{B}).
\]

\begin{table}[H]
\centering
\caption{Proof size (bytes) at representative sizes.
KZG: one $\mathbb{G}_1$ element ($O(1)$);
MKZG: $\mu$ elements ($O(\log n)$);
Bulletproofs: $2\log_2 n + 1$ elements ($O(\log n)$);
Dory: $\sigma$ rounds of $\mathbb{G}_T$, $\mathbb{G}_1$, $\mathbb{G}_2$ elements ($O(\log n)$).}
\label{tab:proofsize-all}
\adjustbox{max width=\textwidth}{%
\begin{tabular}{rlrrrrrrrrrrrr}
\toprule
& & \multicolumn{3}{c}{\textbf{KZG}}
  & \multicolumn{3}{c}{\textbf{MKZG}}
  & \multicolumn{3}{c}{\textbf{Bulletproofs}}
  & \multicolumn{3}{c}{\textbf{Dory}} \\
\cmidrule(lr){3-5}\cmidrule(lr){6-8}\cmidrule(lr){9-11}\cmidrule(lr){12-14}
$n$ & & \textbf{381} & \textbf{254} & \textbf{377}
    & \textbf{381} & \textbf{254} & \textbf{377}
    & \textbf{381} & \textbf{254} & \textbf{377}
    & \textbf{381} & \textbf{254} & \textbf{377} \\
\midrule
4    & & 48   & 32   & 48   & 96   & 64   & 96   & 552   & 392   & 552   & 5\,232  & 3\,488  & 5\,232  \\
16   & & 48   & 32   & 48   & 192  & 128  & 192  & 968   & 680   & 968   & 9\,120  & 6\,080  & 9\,120  \\
64   & & 48   & 32   & 48   & 288  & 192  & 288  & 1\,384 & 968  & 1\,384 & 13\,008 & 8\,672  & 13\,008 \\
256  & & 48   & 32   & 48   & 384  & 256  & 384  & 1\,800 & 1\,256 & 1\,800 & 16\,896 & 11\,264 & 16\,896 \\
1024 & & 48   & 32   & 48   & 480  & 320  & 480  & 2\,216 & 1\,544 & 2\,216 & 20\,784 & 13\,856 & 20\,784 \\
\bottomrule
\end{tabular}}%
\end{table}

\textbf{KZG} produces a constant 48\,B (BLS12-381/377) or 32\,B
(BN254) proof regardless of degree: a single $\mathbb{G}_1$ quotient
commitment checked via one pairing identity.

\textbf{MKZG} adds one $\mathbb{G}_1$ element per variable, giving
$\mu \times 48$\,B on BLS12-381 (96\,B at $\mu=2$, 480\,B at $\mu=10$,
672\,B at $\mu=14$). Still compact, though logarithmically growing.

\textbf{Bulletproofs} sends $2\log_2 n$ group elements ($V_L, V_R$
per round) plus a final element and scalar: 552\,B at $n=4$ growing
to 2\,216\,B at $n=1024$ on BLS12-381, which is $46\times$ larger than
KZG.

\textbf{Dory} is the largest by a wide margin. Each of the
$\sigma = 5$ fold rounds contributes six $\mathbb{G}_T$ elements
(576\,B each on BLS12-381, vs.\ 48\,B for $\mathbb{G}_1$), giving
20\,784\,B at $n=1024$, which is $433\times$ larger than KZG and
$9.4\times$ larger than Bulletproofs. BLS12-381 and BLS12-377 proofs
are identical (same $\mathbb{F}_{p^{12}}$ extension); BN254 is
$33\%$ smaller throughout.

% ============================================================
\section{Impact of Pippenger's MSM Optimisation}
\label{sec:pippenger-comparison}
% ============================================================

Pippenger's multi-scalar multiplication algorithm replaces the naive
$n$ individual scalar multiplications with a batched computation that
has empirical complexity roughly $O(n / \log n)$. This section
quantifies the impact of this optimisation on both KZG and
Bulletproofs, where the naive baseline is directly available.

\subsection{Impact on KZG}

Table~\ref{tab:pip-kzg} shows the commit and prove times for KZG with
and without Pippenger, and the speedup ratio.

\begin{table}[H]
\centering
\caption{Pippenger speedup for KZG commit and prove time on BLS12-381.}
\label{tab:pip-kzg}
\adjustbox{max width=\textwidth}{%
\begin{tabular}{rlllllll}
\toprule
& & \multicolumn{3}{c}{\textbf{Commit (ms)}}
  & \multicolumn{3}{c}{\textbf{Prove (ms)}} \\
\cmidrule(lr){3-5}\cmidrule(lr){6-8}
$n$ & & \textbf{Naive} & \textbf{Pip} & \textbf{Speedup}
    & \textbf{Naive} & \textbf{Pip} & \textbf{Speedup} \\
\midrule
4    & & 1.25  & 0.83  & $1.5\times$ & 1.04  & 0.67  & $1.5\times$ \\
16   & & 3.59  & 1.68  & $2.1\times$ & 3.70  & 1.66  & $2.2\times$ \\
64   & & 14.86 & 4.32  & $3.4\times$ & 15.42 & 4.46  & $3.5\times$ \\
256  & & 56.46 & 12.86 & $4.4\times$ & 56.04 & 12.19 & $4.6\times$ \\
1024 & & 240.78 & 40.37 & $6.0\times$ & 239.35 & 42.10 & $5.7\times$ \\
\bottomrule
\end{tabular}}%
\end{table}

The speedup grows monotonically with $n$: from $1.5\times$ at $n = 4$
to $6.0\times$ at $n = 1024$. This is precisely the expected behaviour
of Pippenger's algorithm, whose advantage over naive MSM grows as
$O(\log n)$. At $n = 4$ there are too few scalars for bucketing to
help, and the overhead of Pippenger's bucket sort marginally reduces
the gap. At $n = 1024$, the $6\times$ speedup is consistent with the
theoretical prediction of $n / \log_b(n) \approx 1024 / \log_2(1024)
= 102$ operations with window size $b \approx 12$, though in practice
memory effects and implementation overhead reduce this to $6\times$.


\begin{table}[H]
\centering
\caption{Pippenger speedup for KZG commit at $n = 1024$ across curves.}
\label{tab:pip-kzg-curves}
\begin{tabular}{lrrrr}
\toprule
\textbf{Curve} & \textbf{Naive (ms)} & \textbf{Pip (ms)} & \textbf{Speedup} \\
\midrule
BLS12-381 & 240.78 & 40.37 & $6.0\times$ \\
BN254     & 123.56 & 20.88 & $5.9\times$ \\
BLS12-377 & 233.76 & 42.00 & $5.6\times$ \\
\bottomrule
\end{tabular}
\end{table}

The speedup is consistent across curves: all three achieve
$5.6$--$6.0\times$ at $n = 1024$. The slightly lower speedup on
BLS12-377 may be due to the larger field modulus (377-bit vs.\
381-bit) interacting with the bucket reduction step of Pippenger's
algorithm.

\subsection{Impact on Bulletproofs}

For Bulletproofs, Pippenger's algorithm most directly benefits the
commitment step, which is a single size-$n$ MSM identical in structure
to KZG commit.

\begin{table}[H]
\centering
\caption{Pippenger speedup for Bulletproofs commit on BLS12-381.}
\label{tab:pip-bp-381}
\begin{tabular}{rllll}
\toprule
$n$ & \textbf{Naive (ms)} & \textbf{Pip (ms)} & \textbf{Speedup} \\
\midrule
16   & 3.90   & 1.42  & $2.8\times$ \\
64   & 18.48  & 3.92  & $4.7\times$ \\
256  & 55.79  & 12.98 & $4.3\times$ \\
1024 & 227.10 & 39.05 & $5.8\times$ \\
\bottomrule
\end{tabular}
\end{table}


\begin{table}[H]
\centering
\caption{Pippenger speedup for Bulletproofs commit at $n = 1024$ across curves.}
\label{tab:pip-bp-curves}
\begin{tabular}{lrrrr}
\toprule
\textbf{Curve} & \textbf{Naive (ms)} & \textbf{Pip (ms)} & \textbf{Speedup} \\
\midrule
BLS12-381 & 227.10 & 39.05 & $5.8\times$ \\
BN254     & 155.18 & 15.41 & $10.1\times$ \\
BLS12-377 & 248.82 & 39.98 & $6.2\times$ \\
\bottomrule
\end{tabular}
\end{table}

The BN254 commit speedup of $10.1\times$ at $n = 1024$ is notably
larger than for BLS12-381/377. This is because BN254's 254-bit scalars
allow Pippenger to use larger effective window sizes for the same
memory budget, and BN254's faster field arithmetic amplifies the
benefit of each bucket operation. This is the largest single Pippenger
speedup observed in all benchmarks.


% ============================================================
\section{Curve Comparison}
\label{sec:curve-comparison}
% ============================================================

The three curves differ in field size (254, 381, 377 bits) and pairing
efficiency. BN254 has a somewhat lower security margin under recent
attacks; BLS12-381 and BLS12-377 offer nominally 128-bit security.
\\\\
\textbf{BN254 vs.\ BLS12-381:}

\begin{table}[H]
\centering
\caption{BN254 speedup over BLS12-381 at $n = 1024$
(value $> 1$ means BN254 is faster).}
\label{tab:curve-ratio}
\begin{tabular}{lrrrr}
\toprule
\textbf{Metric} & \textbf{KZG} & \textbf{MKZG} & \textbf{BP} & \textbf{Dory} \\
\midrule
Setup    & $1.89\times$ & $1.82\times$ & $8.44\times$ & $1.93\times$ \\
Commit   & $1.93\times$ & $1.97\times$ & $2.53\times$ & $1.85\times$ \\
Prove    & $1.79\times$ & $1.77\times$ & $1.56\times$ & $1.53\times$ \\
Verify   & $1.51\times$ & $1.37\times$ & $1.51\times$ & $1.55\times$ \\
Proof    & $1.50\times$ & $1.50\times$ & $1.44\times$ & $1.50\times$ \\
\bottomrule
\end{tabular}
\end{table}

BN254 is $1.5$--$1.9\times$ faster for pairing-heavy operations and
$\approx 2\times$ faster for commit/setup across KZG, MKZG, and Dory.
The outlier is Bulletproofs setup ($8.4\times$): BN254 affine
coordinates are 64 bytes vs.\ 96 bytes for BLS12-381, halving
point-hashing overhead. Proof sizes are uniformly $1.5\times$ smaller
on BN254 (32-byte vs.\ 48-byte $\mathbb{G}_1$ elements), a purely
geometric $3:2$ ratio across all schemes.
\\\\
\textbf{BLS12-381 vs.\ BLS12-377:}

\begin{table}[H]
\centering
\caption{BLS12-381 speedup over BLS12-377 at $n = 1024$.}
\label{tab:curve-ratio-377}
\begin{tabular}{lrrrr}
\toprule
\textbf{Metric} & \textbf{KZG} & \textbf{MKZG} & \textbf{BP} & \textbf{Dory} \\
\midrule
Setup  & $1.13\times$ & $1.01\times$ & $0.97\times$ & $1.14\times$ \\
Commit & $1.04\times$ & $0.99\times$ & $1.02\times$ & $1.12\times$ \\
Prove  & $1.10\times$ & $1.03\times$ & $0.89\times$ & $1.17\times$ \\
Verify & $1.22\times$ & $1.22\times$ & $1.01\times$ & $1.17\times$ \\
\bottomrule
\end{tabular}
\end{table}

The two curves are remarkably close. BLS12-381 is $1.01$--$1.22\times$
faster overall; the largest gap is verification ($1.22\times$ for KZG
and MKZG) because BLS12-381 was optimised for pairing computation,
while BLS12-377 was designed as the inner curve for the ZEXE recursive
SNARK stack. For MSM-only operations, MKZG setup and commit differ by
at most $1.01\times$, well within noise.

Notably, Bulletproofs reverses the trend: BLS12-377 is faster in both
setup ($0.97\times$: 150.0\,ms vs.\ 154.9\,ms) and prove ($0.89\times$:
582.1\,ms vs.\ 650.4\,ms), since Bulletproofs uses only
$\mathbb{G}_1$ scalar multiplications and BLS12-377's 377-bit field
gives marginally more efficient MSM arithmetic without pairings.
Proof sizes are identical (both use 48-byte $\mathbb{G}_1$ points).\\\\
\textbf{Curve Selection Guidance:}

Choose \textbf{BN254} for best latency and smallest proofs when
128-bit security (under current attacks) is acceptable.
Choose \textbf{BLS12-381} for a higher security margin or Ethereum
precompile compatibility, at an $\approx 1.9\times$ performance cost
over BN254. Choose \textbf{BLS12-377} when recursive SNARK
composability is required. Its performance is near-identical to
BLS12-381.




% ============================================================
\section{Overall Scheme Comparison}
\label{sec:overall}
% ============================================================

Table~\ref{tab:overall-n1024} provides a side-by-side summary of all
benchmarked metrics at the largest tested size ($n = 1024$, or
$\mu = 10$ for MKZG) on BLS12-381, using Pippenger for KZG and
Bulletproofs.

\begin{table}[H]
\centering
\caption{Full performance summary at $n = 1024$ on BLS12-381.}
\label{tab:overall-n1024}
\begin{tabular}{lrrrr}
\toprule
\textbf{Metric} & \textbf{KZG} & \textbf{MKZG} & \textbf{Bulletproofs} & \textbf{Dory} \\
\midrule
Setup (ms)      & 1027.5 & 235.9  & 154.9   & 111.0  \\
Commit (ms)     & 40.4   & 39.5   & 39.1    & 91.0   \\
Prove (ms)      & 42.1   & 263.4  & 650.4   & 264.3  \\
Verify (ms)     & 3.8    & 21.9   & 496.3   & 95.8   \\
Proof (bytes)   & 48     & 480    & 2\,216  & 20\,784 \\
Trusted setup   & Yes    & Yes    & No      & No     \\
Polynomial type & Univar.& Multilin.& Vector & Multilin. \\
\bottomrule
\end{tabular}
\end{table}

\begin{table}[H]
\centering
\caption{Full performance summary at $n = 1024$ on BN254.}
\label{tab:overall-n1024-bn254}
\begin{tabular}{lrrrr}
\toprule
\textbf{Metric} & \textbf{KZG} & \textbf{MKZG} & \textbf{Bulletproofs} & \textbf{Dory} \\
\midrule
Setup (ms)      & 544.9  & 129.6  & 18.4    & 57.6   \\
Commit (ms)     & 20.9   & 20.1   & 15.4    & 49.3   \\
Prove (ms)      & 23.5   & 149.2  & 415.7   & 173.1  \\
Verify (ms)     & 2.5    & 16.0   & 329.6   & 61.9   \\
Proof (bytes)   & 32     & 320    & 1\,544  & 13\,856 \\
Trusted setup   & Yes    & Yes    & No      & No     \\
\bottomrule
\end{tabular}
\end{table}

Figures~\ref{fig:combined-bls12381} and~\ref{fig:combined-bn254} plot
setup, commit, prove, and verify times on a log scale for all four
schemes across the full size range $n = 4$ to $n = 16{,}384$,
complementing the tabular summaries above.

\begin{figure}[H]
\centering
\includegraphics[width=\textwidth]{"../Final Graphs/plots/BLS12381_combined"}
\caption{All four schemes on BLS12-381: setup, commit, prove, and verify
time (ms, log scale) vs.\ polynomial size $n$.}
\label{fig:combined-bls12381}
\end{figure}

\begin{figure}[H]
\centering
\includegraphics[width=\textwidth]{"../Final Graphs/plots/BN254_combined"}
\caption{All four schemes on BN254: setup, commit, prove, and verify
time (ms, log scale) vs.\ polynomial size $n$.}
\label{fig:combined-bn254}
\end{figure}

% ============================================================
\section{Practical Recommendations}
\label{sec:recommendations}
% ============================================================

\begin{itemize}
    \item \textbf{Univariate SNARKs:} Use \textbf{KZG on BN254} with
          Pippenger. It gives 32\,B proofs, 2.5\,ms verification, and
          a prover $6$--$15\times$ faster than all alternatives. The
          trusted setup is a one-time multi-party ceremony.

    \item \textbf{Multilinear commitments (Spartan, HyperPlonk):}
          Use \textbf{MKZG on BN254} if a trusted setup is acceptable
          (16\,ms verify, 320\,B proof at $\mu = 10$). Otherwise use
          \textbf{Dory}: $O(\log n)$ verification without a trusted
          setup. Avoid Bulletproofs above $n = 64$; its $O(n)$ verify
          becomes prohibitive.

    \item \textbf{On-chain / proof-size critical:} Use \textbf{KZG on
          BN254}. Its 32\,B proof is $46\times$ smaller than
          Bulletproofs and $433\times$ smaller than Dory; BN254 is
          directly supported by Ethereum precompiles.

    \item \textbf{Pairing-free environments:} Use \textbf{Bulletproofs};
          it requires only $\mathbb{G}_1$ operations, unlike Dory which
          needs $\mathbb{G}_1$, $\mathbb{G}_2$, and $\mathbb{G}_T$.

    \item \textbf{Recursive SNARK stacks:} Use \textbf{BLS12-377}
          (inner curve in ZEXE/BLS12-378). Use \textbf{BLS12-381} for
          higher security or Ethereum compatibility; \textbf{BN254}
          for maximum performance.
\end{itemize}

% ============================================================
\section{Summary}
\label{sec:chapter7-summary}
% ============================================================

This chapter has benchmarked and compared four polynomial commitment
schemes across six metrics, three curves, and nine polynomial sizes.
The main findings are:

\begin{enumerate}
    \item \textbf{KZG dominates on prover efficiency and proof size.}
          With Pippenger's MSM, KZG proves in 42\,ms and produces
          48\,byte proofs on BLS12-381 at $n = 1024$. Its verifier is
          constant-time (3.8\,ms), making it optimal for deployed
          SNARK back-ends. The required trusted setup is the only
          practical limitation.

    \item \textbf{MKZG extends KZG to multilinear polynomials}. It is the correct choice for sumcheck-based proof systems where the polynomial structure is inherently multilinear.

    \item \textbf{Dory is the best transparent scheme for large
          polynomials.} It achieves $O(\log n)$ verification without a
          trusted setup, at the cost of large proofs (20\,KB at
          $n = 1024$) and higher proving latency than KZG.

    \item \textbf{Bulletproofs are limited by $O(n)$ verification.}
          The 496\,ms verify time at $n = 1024$ on BLS12-381 makes
          Bulletproofs unsuitable for large polynomials. They remain
          useful for small $n$ in pairing-free environments.

    \item \textbf{Pippenger's algorithm delivers $5$--$6\times$
          commit/prove speedup} for KZG and Bulletproofs wherever a
          single large MSM dominates. It provides no benefit for
          operations that are inherently sequential (KZG setup) or
          composed of many small MSMs (Bulletproofs verify).

    \item \textbf{BN254 is $\approx 1.9\times$ faster than BLS12-381}
          for all schemes and all metrics, with identical algorithmic
          structure. BLS12-377 and BLS12-381 are nearly equal in
          performance. Proof sizes on BN254 are $1.5\times$ smaller
          due to the smaller field.

    \item \textbf{The verification hierarchy} (fastest to slowest at
          $n = 1024$, BLS12-381) is:
          KZG (3.8\,ms) $\ll$ MKZG (21.9\,ms) $<$ Dory (95.8\,ms)
          $\ll$ Bulletproofs (496\,ms).
          This ordering is stable across all three curves and grows
          with $n$ in a way that further separates the linear verifier
          (Bulletproofs) from the logarithmic and constant ones.
\end{enumerate}

No single scheme dominates on all dimensions simultaneously. The
correct choice depends on whether a trusted setup is acceptable, what
polynomial type the application uses, how proof size is weighted
against verifier latency, and which elliptic curve is mandated by the
deployment environment. The benchmarks in this chapter provide
concrete, curve-specific numbers to inform this decision for any
target application.
